\documentclass[11pt]{amsart}
\usepackage[T1]{fontenc}
\usepackage[utf8]{inputenc}
\usepackage[french]{babel}
\usepackage{amssymb}
\usepackage[margin=1.05in]{geometry}
\usepackage[hidelinks]{hyperref}
\newcommand{\stamptime}{\the\numexpr\time/60\relax:\ifnum\numexpr\time-60*(\time/60)\relax<10 0\fi\the\numexpr\time-60*(\time/60)\relax}
\AddToHook{shipout/foreground}{\put(0.5\paperwidth,-\paperheight+0.55cm){\makebox[0pt]{\tiny\ttfamily unit/doc/\jobname.tex \ --- \ compiled \today, \stamptime}}}
\newtheorem{proposition}{Proposition}
\newtheorem{lemma}[proposition]{Lemme}
\newtheorem{definition}[proposition]{Définition}
\theoremstyle{remark}
\newtheorem{remark}[proposition]{Remarque}
\DeclareMathOperator{\im}{im}
\newcommand{\La}{\Lambda}
\newcommand{\Rees}{\mathcal R}
\title{ACL \S4 par \v{C}ech : preuve informelle des parties \`a formaliser}
\date{26 octobre 2026}
\begin{document}
\maketitle

\noindent\textbf{Statut.}
Cette note complète \texttt{finite-lifting.tex} et
\texttt{finite-cohomology.tex}. Elle donne, sous une forme découpée pour
Unit, les preuves des étapes du jalon~3 de \texttt{acl-zmt-handoff.txt} :
le modèle explicite du complexe de Čech de $B_h$, la restriction à la
carte $Z$, la donnée de finitude (F1) à partir d'une donnée d'élimination,
et la donnée (F2) à partir d'un lemme de syzygies. Les étapes déjà
vérifiées dans \texttt{sag/acl-zmt} sont seulement citées
(\S\ref{sec:verifie}). Deux entrées restent des hypothèses explicites :
la donnée d'élimination (E) et la donnée de syzygies (S). Dans l'exemple
d'ACL, (S0) est triviale ($D_0=1$) et (S1) a une preuve courte
(remarque~\ref{rem:S1acl}).

\section{Cadre}

$A$ est un anneau commutatif, $a,b\in A$, $I=(a,b)$. $S_0=A[X,Y,Z]$ est
gradué par le degré total. $P,Q\in S_0$ sont homogènes de degrés $d,e\ge1$,
homogénéisés de $p=x+af+bg$ et $q=y+au+bv$ :
\[
 P=XZ^{d-1}+aF+bG,\qquad Q=YZ^{e-1}+aU+bV .
\]
On note $B=A[x,y]/(p,q)$ (ou tout $A$-algèbre $L$ avec $x,y\in L$,
$p(x,y)=q(x,y)=0$ : tous les énoncés sont de cette forme).

$\La$ désigne les polynômes de Laurent en $X,Y,Z$ à coefficients dans $A$,
donnés par des listes finies de termes $cX^m$, $m\in\mathbf Z^3$ ; leur
\emph{sens} est la fonction coefficient $m\mapsto\sum_{(c,m')\,:\,m'=m}c$
(\texttt{Laurent3}). Pour une partie non vide $E\subseteq\{X,Y,Z\}$,
$\La_E(k)$ est l'ensemble des éléments de degré $k$ dont les exposants
des variables hors de $E$ sont $\ge0$.

\section{Ce qui est vérifié}\label{sec:verifie}
Correction finie ; contractions $\pi,\sigma,\tau$ et génération de
$H^1$ (\texttt{CechH1Finite}, sur un anneau quelconque $D$) ; cochaîne
fixe $Z^{-1}=E^nZ^{-1}+Pu_n$ (\texttt{FixedCochain}) ; relation graduée
par noethérianité codée et relèvement $h$ (\texttt{GradedDependence},
\texttt{ReesLifting}) ; équations unitaires par Cayley--Hamilton à partir
de générateurs et de constantes de structure (\texttt{IntegralSections}) ;
le théorème conditionnel \texttt{ACLZMT}.

\section{Le modèle de Čech de $B_h$}\label{sec:modele}

\begin{definition}
$K^0(k)=\La_X(k)\times\La_Y(k)\times\La_Z(k)$,
$K^1(k)=\La_{XY}(k)\times\La_{XZ}(k)\times\La_{YZ}(k)$,
$K^2(k)=\La_{XYZ}(k)$, avec $\delta_0,\delta_1$ de \texttt{Cech3}.
Deux $j$-cochaînes sont \emph{congrues}, $f\equiv f'$, si
$f-f'=Pu+Qv$ composante par composante, avec $u$ (resp.\ $v$) une
$j$-cochaîne de degré $k-d$ (resp.\ $k-e$) vérifiant les mêmes
conditions de support. On note $N^j(k)$ ces différences.
\end{definition}

On ne prouve pas que ce complexe calcule la cohomologie du faisceau
$\mathcal O_W(k)$ : le théorème final n'en a pas besoin. Ce qui est
utilisé est la liste suivante, chacune élémentaire.

\begin{lemma}\label{lem:modele}
\begin{enumerate}
\item $N^j(k)$ est un sous-module, stable par $\delta$ :
  $\delta(Pu+Qv)=P\,\delta u+Q\,\delta v$ ; et par la multiplication
  par $X,Y,Z$ (qui envoie $N^j(k)$ dans $N^j(k+1)$).
\item $K^0(0)$ est un anneau composante par composante ; $N^0(0)$ en est
  un idéal ; si $\delta_0c,\delta_0c'\in N^1(0)$ alors
  $\delta_0(cc')\in N^1(0)$, car
  $(cc')_j-(cc')_i=c_j(c'_j-c'_i)+(c_j-c_i)c'_i$.
\item Restriction. Pour $f\in\La_Z(k)$ (exposants de $X,Y$ positifs),
  $\mathrm{res}_k(f)=\sum c\,x^{m_X}y^{m_Y}$ sur les termes $cX^m$ de $f$.
  On a $\mathrm{res}_{k+k'}(fg)=\mathrm{res}_k(f)\,\mathrm{res}_{k'}(g)$,
  $\mathrm{res}_d(P)=p(x,y)=0$, $\mathrm{res}_e(Q)=0$, donc
  $\mathrm{res}_k$ s'annule sur $P\La_Z(k-d)+Q\La_Z(k-e)$.
  Enfin $\mathrm{res}_0(Zw)=\mathrm{res}_{-1}(w)$,
  $\mathrm{res}_0(Xw)=x\,\mathrm{res}_{-1}(w)$,
  $\mathrm{res}_0(Yw)=y\,\mathrm{res}_{-1}(w)$ et
  $\mathrm{res}_0(Z\cdot Z^{-1})=1$.
\item $\mathrm{res}_k(f)$ ne dépend que du sens de $f$ (lemme suivant).
\end{enumerate}
\end{lemma}

\begin{lemma}[regroupement]\label{lem:regroupement}
Soit $\varphi:\mathbf Z^3\to L$ quelconque. Si la fonction coefficient
d'une liste $\ell$ est nulle, alors $\sum_{(c,m)\in\ell}c\,\varphi(m)=0$.
\end{lemma}
\begin{proof}
Récurrence sur une borne de la longueur. Si $\ell=(c,m)::\ell'$,
séparer $\ell$ en $\ell_{=m}$ (termes d'exposant $m$, comparaison
d'entiers) et $\ell_{\ne m}$. La somme sur $\ell_{=m}$ vaut
$(\text{coef}_\ell(m))\varphi(m)=0$ ; la fonction coefficient de
$\ell_{\ne m}$ est nulle et $\ell_{\ne m}$ est strictement plus courte.
\end{proof}

\begin{remark}[représentation]
Le théorème \texttt{ACLZMT} prend des modules à égalité $\mathrm{Id}$.
Deux réalisations sont possibles : des modules quotients (ensembles
quotients des listes) ou des modules de représentants munis de la
congruence $\equiv$. La seconde est recommandée : elle garde explicites
les représentants utilisés aux sections~\ref{sec:F1}--\ref{sec:F2}, alors
qu'un quotient ne fournit ses représentants que sous troncation. Il
faut pour cela remplacer, dans \texttt{FiniteCorrection},
\texttt{ReesLifting} et \texttt{ACLZMT}, les égalités dans $K^1$ et $C^1$
par l'appartenance à $N^1$ ; pour $K^1$ c'est un transport le long d'une
égalité exacte. Mais il faut aussi affaiblir la génération en degré zéro :
(F2) ne donne que $c\equiv a_0+\sum_k\mu_kc_{t_k}\pmod{N^0}$, alors que
\texttt{ACLZMT.gen2} et \texttt{one-comb} demandent une égalité exacte
dans $C^0$. On remplace donc ces égalités par des congruences modulo $N^0$,
et les preuves de \texttt{ACLZMT} (\texttt{res-comb}, \texttt{table},
\texttt{monic}, \texttt{one}) utilisent que la restriction annule $N^0$
(lemme~\ref{lem:modele}(3)) ; ce n'est pas un simple transport. Il faut aussi remplacer
l'hypothèse générale \texttt{gen1} par sa seule instance utilisée, les
classes $\beta_n=[\delta_0z_*]$ (ce qui affaiblit l'hypothèse).
Les éléments de Laurent produits par \texttt{FixedCochain} sont des
fonctions ; pour la restriction il faut leurs listes : les opérateurs
$E=a(-\theta F)+b(-\theta G)$ s'écrivent sur les listes (produit de
listes), avec $\text{coef}(g\cdot\ell)=g\cdot\text{coef}(\ell)$
(\texttt{mulP}), ce qui demande seulement l'arithmétique de $\mathbf Z$
déjà utilisée par \texttt{mulP-comm}.
\end{remark}

\section{La cochaîne fixe dans le modèle}

\begin{proposition}\label{prop:fixe}
$z_*=(0,0,Z^{-1})\in K^0(-1)$ vérifie $\mathrm{res}_0(Zz_*)=1$ et, pour
tout $n$, $\delta_0z_*\equiv w_n$ avec $w_n\in I^nK^1(-1)$, par un arbre
explicite dont les feuilles vérifient les conditions de support et sont
de degré $-1$.
\end{proposition}
\begin{proof}
$\delta_0z_*=(0,Z^{-1},Z^{-1})$. Sur $XZ$,
$Z^{-1}=E^nZ^{-1}+Pu_n$ ; sur $YZ$, $Z^{-1}=E'^nZ^{-1}+Qu'_n$
(\texttt{FixedCochain.Chart.fixed}). On pose
$w_n=(0,E^nZ^{-1},E'^nZ^{-1})$. Les deux arbres ont la même profondeur
$n$ ; on les combine composante par composante : un noeud
$(0,f,g)=a(0,f_1,g_1)+b(0,f_2,g_2)$, avec la composante $XY$ nulle.
\end{proof}

\section{La donnée (F1) à partir d'une élimination}\label{sec:F1}

Soit $D=A[U,V]$, bigradué : degré projectif en $X,Y,Z$, degré de Rees en
$U,V$. Soit $\La^D$ les polynômes de Laurent sur $D$ et
$\phi:\La^D\to\La[T]$ le morphisme $U\mapsto aT$, $V\mapsto bT$,
coefficient par coefficient en $X,Y,Z$. Si $f$ est de degré de Rees $n$,
$\phi(f)=\bar fT^n$ et $\bar f\in I^n\La$, par un arbre explicite
(un monôme $U^\alpha V^\beta$ donne $a^\alpha b^\beta$).

\begin{definition}[donnée d'élimination (E)]
Une liste $g_1,\dots,g_r\in D[X,Y,Z]$, $g_j$ bihomogène de bidegré
$(d_j,e_j)$, avec
\begin{itemize}
\item[(E1)] $\bar g_j=P\alpha_j+Q\beta_j$ dans $S_0$ (témoins explicites) ;
\item[(E2)] pour tout $F\in D[X,Y,Z]$ bihomogène tel que
 $\bar F=P\alpha+Q\beta$ avec $\alpha,\beta\in S_0$, une écriture
 $F=\sum_jg_jh_j$, $h_j\in D[X,Y,Z]$.
\end{itemize}
\end{definition}

(E) dit que $(g_j)$ engendre $J=\ker(D[X,Y,Z]\to B_h[T])$. Elle est
fournie par une base de Gröbner de $(P,Q,U-aT,V-bT)$ pour un ordre
éliminant $T$ (Coquand--Persson, théorème 11) : (E2) est alors la
division par la base. C'est l'entrée principale qui reste à formaliser.

\begin{lemma}[homogénéité]\label{lem:homog}
Si $F=\sum_jg_jh_j$ avec $F$ et les $g_j$ bihomogènes, alors
$F=\sum_jg_jh'_j$ où $h'_j$ est la composante de $h_j$ de bidegré
$\deg F-\deg g_j$ ; dans $D[X,Y,Z]$ elle est nulle si l'un des deux
degrés est négatif. Même énoncé dans $\La^D$, où seul un degré de
\emph{Rees} négatif force l'annulation : le degré projectif d'un élément
de Laurent peut être négatif ($1=X\cdot X^{-1}$).
\end{lemma}
\begin{proof}
Prendre la composante de bidegré $\deg F$ des deux membres : la
projection est linéaire et commute à la multiplication par un élément
bihomogène, en décalant le degré. Coefficient par coefficient, c'est
\texttt{hom-mul} de \texttt{Laurent3} plus une projection explicite.
\end{proof}

\begin{lemma}[dénominateurs]\label{lem:denom}
Pour toute liste $f\in\La_E$ il existe $K$ tel que $\mu f$ soit un
polynôme, $\mu=\prod_{V\in E}V^K$ : prendre $K$ plus grand que l'opposé
de tous les exposants négatifs des termes de $f$ (support fini).
\end{lemma}

\begin{proposition}[F1 pour les classes $\beta_n$]\label{prop:F1}
Sous (E), posons $\mathcal T=\{(j,t)\,:\,t\in\mathrm{gens}(d_j)\}$ et,
pour $(j,t)\in\mathcal T$, $m_{j,t}=\sigma(\bar g_jX^{\mathrm{expo}\,t})
\in K^1(-1)$, de degré de Rees $e_j$. Pour tout $n$ on construit
$\rho_n\in A^{\mathcal T}$, avec $\rho_{n,(j,t)}\in I^{n-e_j}$ et
$\rho_{n,(j,t)}=0$ si $n<e_j$, et $v_n\in I^nK^0(-1)$ tels que
\[
 \delta_0z_*\equiv\sum_{(j,t)}\rho_{n,(j,t)}\,m_{j,t}+\delta_0v_n .
\]
\end{proposition}
\begin{proof}
\emph{Relèvement.} Soit l'arbre de $w_n$ (proposition~\ref{prop:fixe}).
On pose $\mathrm{Lift}(\text{feuille }\ell)=\ell$ et
$\mathrm{Lift}(am_1+bm_2)=U\,\mathrm{Lift}(m_1)+V\,\mathrm{Lift}(m_2)$.
Alors $W=\mathrm{Lift}(w_n)$ est une 1-cochaîne de $\La^D$, de degré de
Rees $n$, vérifiant les supports, et $\phi(W)=w_nT^n$.

\emph{Cocycle modulo $J$.} $\delta_1w_n=\delta_1(\delta_0z_*-(\delta_0z_*-w_n))
=-\delta_1(Pu+Qv)=Pu'+Qv'$ (car $\delta_1\delta_0=0$), avec $u',v'$
explicites dans $\La_{XYZ}$. Donc $\phi(\delta_1W)=T^n(Pu'+Qv')$. Par le
lemme~\ref{lem:denom}, choisir $\mu=(XYZ)^K$ avec $\mu\,\delta_1W$,
$\mu u'$, $\mu v'$ polynomiaux ; alors
$\overline{\mu\delta_1W}=P(\mu u')+Q(\mu v')$ et (E2) donne
$\mu\,\delta_1W=\sum_jg_jh_j$. Par le lemme~\ref{lem:homog}, les $h_j$
sont bihomogènes ; $q_j=h_j/\mu\in\La^D_{XYZ}$ est de degré projectif
$-1-d_j$ et de degré de Rees $n-e_j$, et $\delta_1W=\sum_jg_jq_j$.

\emph{Génération sur $D$.} \texttt{CechH1Finite.Reduce.generation} sur
l'anneau $D$ donne
\[
 W=\sum_jg_j\sigma(q_j)
   +\sum_{(j,t)}q_j(\mathrm{expo}\,t)\,\sigma(g_jX^{\mathrm{expo}\,t})
   +\delta_0c ,\qquad c=\tau(\cdots).
\]
$\sigma,\tau,\delta$ agissent coefficient par coefficient en $X,Y,Z$ et
sont $D$-linéaires : ils préservent le degré de Rees et commutent à
$\phi$. Le coefficient $\lambda=q_j(\mathrm{expo}\,t)\in D$ est de degré
de Rees $n-e_j$, donc $\phi(\lambda)=\rho T^{n-e_j}$ avec
$\rho\in I^{n-e_j}$, et $\lambda=0$ si $n<e_j$.

\emph{Retour.} Appliquer $\phi$ et prendre le coefficient de $T^n$ :
$\phi(g_j\sigma(q_j))$ donne $\bar g_j\,\overline{\sigma(q_j)}
=P(\alpha_j\overline{\sigma(q_j)})+Q(\beta_j\overline{\sigma(q_j)})
\equiv0$ par (E1) ; le terme suivant donne $\rho_{n,(j,t)}m_{j,t}$ ; le
dernier donne $\delta_0v_n$, $v_n=\bar c\in I^nK^0(-1)$. Comme
$\delta_0z_*\equiv w_n$, la congruence suit.
\end{proof}

C'est exactement l'hypothèse de \texttt{ReesLifting}, avec $G$ l'application
$\rho\mapsto\sum\rho_{(j,t)}m_{j,t}$ et $e_{(j,t)}=e_j$ ; le nombre de
générateurs est $\sum_j\binom{d_j}{2}$. Seules les données jusqu'à
l'indice $N$ fourni par la noethérianité sont construites.

\section{La donnée (F2) à partir des syzygies}\label{sec:F2}

Soit $D_0$ un diviseur commun homogène de $P,Q$, $P=D_0P'$,
$Q=D_0Q'$, $r=\deg D_0$, $\delta=d+e-r$.

\begin{definition}[donnée de syzygies (S)]
\begin{itemize}
\item[(S0)] $D_0f=0$ dans $\La_{XYZ}$ entraîne $f=0$ ;
\item[(S1)] si $P'\alpha+Q'\beta=0$ dans $\La_{XYZ}$, avec $\alpha,\beta$
  homogènes, on construit $w$ avec $\alpha=-Q'w$, $\beta=P'w$.
\end{itemize}
\end{definition}

(S1) se ramène à $S_0$ par le lemme~\ref{lem:denom} : si
$\mu\alpha=-Q'\tilde w$, $\mu\beta=P'\tilde w$, poser $w=\tilde w/\mu$.
Sur $S_0$, c'est la coprimalité de $P',Q'$ (pgcd dans
$k[a,b,X,Y,Z]$, $k$ corps discret).

\begin{remark}[(S1) pour l'exemple d'ACL]\label{rem:S1acl}
$A=k[a,b]$, $P=XZ-aZ^2+b(XY+2X^2)$, $Q=YZ+a(X^2+XY)+b(Y^2-Z^2)$, $D_0=1$.
Dans $\La_{XYZ}$ le coefficient $-Z^2$ de $a$ dans $P$ est inversible :
$\La_{XYZ}/(P)\simeq k[b,X^{\pm1},Y^{\pm1},Z^{\pm1}]$ en substituant
$a=Z^{-2}(XZ+b(XY+2X^2))$. Ce quotient est intègre (anneau de Laurent sur
$k[b]$, terme dominant pour un ordre de groupe sur $\mathbf Z^3$, $k$
discret) et l'image $\bar Q$ de $Q$ est non nulle (un coefficient
explicite). Si $P\alpha+Q\beta=0$, alors $\bar Q\bar\beta=0$, donc
$\beta=Pw$ ; puis $P(\alpha+Qw)=0$ et $P$ est régulier dans $\La_{XYZ}$
(intègre), d'où $\alpha=-Qw$.
\end{remark}

Soit $\mathcal L\subseteq H^2(\mathcal O(-\delta))$ (combinaisons des
monômes tous négatifs de degré $-\delta$) l'ensemble des $t$ tels que
$\pi(Q't)=\pi(P't)=0$. Pour $t\in\mathcal L$ :
\[
 u_1=\sigma(-Q't),\quad u_2=\sigma(P't),\quad c_t=\tau(Pu_1+Qu_2).
\]

\begin{lemma}\label{lem:ct}
$\delta_0c_t=Pu_1+Qu_2$ exactement ; donc $c_t\in K^0(0)$ est fermée
modulo $N$. L'application $t\mapsto c_t$ est $A$-linéaire.
\end{lemma}
\begin{proof}
$\delta_1(Pu_1+Qu_2)=P(-Q't-\pi(-Q't))+Q(P't-\pi(P't))=-PQ't+QP't=0$,
car $PQ'=QP'=D_0P'Q'$ ; les supports sont ceux de $\sigma$ ; on applique
\texttt{contraction}. Linéarité : $\sigma,\tau$ et les produits le sont.
\end{proof}

\begin{proposition}[F2]\label{prop:F2}
Soit $c\in K^0(0)$ avec $\delta_0c=Pu+Qv$ ($u$, $v$ des 1-cochaînes de
degrés $-d,-e$ vérifiant les supports). Alors on construit $a_0\in A$ et
$t\in\mathcal L$ tels que $c\equiv a_0\cdot1+c_t$. Si $t_1,\dots,t_s$
engendrent $\mathcal L$ (avec les coordonnées explicites), alors
$1,c_{t_1},\dots,c_{t_s}$ engendrent les sections fermées.
\end{proposition}
\begin{proof}
$0=\delta_1\delta_0c=P\,\delta_1u+Q\,\delta_1v=D_0(P'\delta_1u+Q'\delta_1v)$.
Par (S0) puis (S1), $\delta_1u=-Q'w$, $\delta_1v=P'w$ avec $w$ de degré
$-\delta$. Décomposer $w=\pi(w)+\delta_1\sigma(w)$ (\texttt{decomp}) et
poser $t=\pi(w)$. Alors
\[
 \delta_1(u+Q'\sigma(w))=-Q't,\qquad \delta_1(v-P'\sigma(w))=P't .
\]
Un $\delta_1$ de 1-cochaîne vérifiant les supports n'a pas de monôme tout
négatif ; donc $\pi(Q't)=\pi(P't)=0$ : $t\in\mathcal L$. Posons
$u''=u+Q'\sigma(w)-u_1$ et $v''=v-P'\sigma(w)-u_2$ : par
\texttt{decomp}, $\delta_1u''=\delta_1v''=0$, d'où
$u''=\delta_0\tau(u'')$, $v''=\delta_0\tau(v'')$
(\texttt{contraction}, tout degré). Le calcul
\[
 \delta_0c=Pu''+Qv''+(Pu_1+Qu_2)+(QP'-PQ')\sigma(w)
          =\delta_0\big(P\tau(u'')+Q\tau(v'')+c_t\big)
\]
montre que $c'=c-c_t-P\tau(u'')-Q\tau(v'')$ vérifie $\delta_0c'=0$
exactement. Ses trois composantes sont alors égales à une même fonction
dont tous les exposants sont $\ge0$ (chaque variable est positive dans
deux des trois cartes), de degré $0$ : $c'=a_0\cdot1$, $a_0=c'(0,0,0)$.
Donc $c\equiv a_0+c_t$. Avec $t=\sum\mu_kt_k$ et la linéarité,
$c\equiv a_0+\sum\mu_kc_{t_k}$.
\end{proof}

Le calcul de générateurs de $\mathcal L$ (noyau d'une matrice à
coefficients dans $A$) est la seconde élimination du handoff (\S8). Dans
l'exemple d'ACL, $d=e=2$, $r=0$ : $H^2(\mathcal O(-2))=0$, donc
$\mathcal L=H^2(\mathcal O(-4))=A^3$, de base les trois monômes, et les
coordonnées sont les valeurs $t(\mathrm{expo})$ (comme
\texttt{pi-gens}). Les éléments $m_k$ de \texttt{ACLExample} sont les
$\mathrm{res}_0(c_{t_k})$.

\section{Assemblage}

\begin{proposition}
Soit $A=k[a,b]$, $k$ corps discret. Sous (E) et (S), avec des générateurs
explicites de $\mathcal L$, il existe sous troncation $s\in B$ avec
$s=1+ar+bt$ et $s,sx,sy$ racines de polynômes unitaires sur $A$.
\end{proposition}
\begin{proof}
Instancier \texttt{ACLZMT} (sous la forme à congruence de la
remarque de la section~\ref{sec:modele}) : le modèle de Čech et ses
champs par le lemme~\ref{lem:modele} ; la cochaîne fixe par la
proposition~\ref{prop:fixe} ; (F1) pour les $\beta_n$ par la
proposition~\ref{prop:F1} ; (F2) par la proposition~\ref{prop:F2},
appliquée aux représentants explicites $Zh,Xh,Yh$ et aux produits
$c_ic_j$ (fermés par le lemme~\ref{lem:modele}) ; ses congruences modulo
$N^0$ deviennent des égalités après restriction. La présentation
noethérienne codée de $k[a,b]$ : \texttt{Field.noeth} pour $k$,
\texttt{Transfer} le long de $k\to A$, \texttt{Generated} pour $a,b$ ;
tout élément de $A$ se relève à un code sous troncation, par induction
sur le quotient qui définit $A[X]$.
\end{proof}


\section{Deux raccourcis : ni base de Gröbner, ni générateurs de $\mathcal L$}\label{sec:raccourcis}

\noindent\textit{À faire vérifier avant formalisation.} Ils remplacent (E)
et les générateurs de $\mathcal L$ ; seule (S) reste une hypothèse.

\subsection*{(F1) avec un seul élément de $J$}
Sur la carte $XYZ$, posons $\alpha=-\theta F$, $\beta=-\theta G$
($\theta=X^{-1}Z^{1-d}$), $\alpha'=-\theta'U$, $\beta'=-\theta'V$
($\theta'=Y^{-1}Z^{1-e}$), de sorte que $a\alpha+b\beta=1-\theta P$ et
$a\alpha'+b\beta'=1-\theta'Q$ (\texttt{FixedCochain.E-eq}). Dans $\La^D$ :
\[
 \tilde\alpha=U\alpha+V\beta,\qquad \tilde\alpha'=U\alpha'+V\beta' ,
 \qquad \phi(\tilde\alpha)=T(1-\theta P),\ \phi(\tilde\alpha')=T(1-\theta'Q).
\]
Le relèvement de l'arbre de $w_n$ (proposition~\ref{prop:F1}) est
$W=(0,\tilde\alpha^nZ^{-1},\tilde\alpha'^nZ^{-1})$, d'où
\[
 \delta_1W=(\tilde\alpha'^n-\tilde\alpha^n)Z^{-1}
 =(\tilde\alpha'-\tilde\alpha)\sum_{i<n}\tilde\alpha'^{\,i}\tilde\alpha^{\,n-1-i}Z^{-1}.
\]
Soit $\mu$ un monôme tel que $g_*=\mu(\tilde\alpha'-\tilde\alpha)$ soit un
polynôme ; $g_*$ est bihomogène, de degré de Rees $1$, et
$\phi(g_*)=T\mu(\theta P-\theta'Q)$, ce qui est (E1) pour $g_*$ avec des
témoins explicites. Avec $q=\mu^{-1}\sum_{i<n}\tilde\alpha'^{\,i}
\tilde\alpha^{\,n-1-i}Z^{-1}$, on a $\delta_1W=g_*q$ exactement : la
preuve de la proposition~\ref{prop:F1} s'applique avec la seule liste
$(g_*)$, sans (E2). Les générateurs sont
$m_t=\sigma(\bar g_*X^{\mathrm{expo}\,t})$, $t\in\mathrm{gens}(\deg g_*)$,
tous de degré de Rees $1$. La finitude ne vient que de la noethérianité
codée de $A$ (\texttt{GradedDependence}) : on n'a besoin que des $\beta_n$,
pas de tout $M$.

\subsection*{(F2) : l'équation unitaire directement}
La preuve de la proposition~\ref{prop:F2} associe à toute section fermée
$c$ un couple $\kappa(c)=(a_0,t)\in A\times H^2(\mathcal O(-\delta))$ avec
$c\equiv a_0+c_t$ ; $H^2(\mathcal O(-\delta))$ est libre de base les
monômes tous négatifs de degré $-\delta$, donc
$\kappa(c)\in A^{1+m}$, $m=\binom{\delta-1}{2}$. On n'utilise pas de
générateurs de $\mathcal L$.

Soit $c$ fermée (ici $Zh$, $Xh$ ou $Yh$) et $c_i=c^i$ (fermées par le
lemme~\ref{lem:modele}). La noethérianité codée de $A$, sous la forme non
graduée de \texttt{GradedDependence} (tous les $e_k=0$, sans $T$), donne
sous troncation $N$ et $\lambda_i\in A$ avec
$\kappa(c_N)=\sum_{i<N}\lambda_i\kappa(c_i)$. Par linéarité de
$t\mapsto c_t$,
\[
 c^N-\sum_{i<N}\lambda_ic^i\equiv\Big(a_0(c_N)-\sum_{i<N}\lambda_ia_0(c_i)\Big)
   +c_{\,t(c_N)-\sum_i\lambda_it(c_i)}=0+c_0=0 ,
\]
c'est-à-dire $c^N\equiv\sum_{i<N}\lambda_ic^i\pmod{N^0}$. La restriction
annule $N^0$ et est multiplicative : $\mathrm{res}(c)$ est racine du
polynôme unitaire $X^N-\sum\lambda_iX^i$. On n'a besoin ni de la linéarité
de $\kappa$ (seulement de celle de $t\mapsto c_t$), ni de
\texttt{IntegralSections}.

\subsection*{Ce qui reste}
La donnée (S) : pour l'exemple d'ACL, la remarque~\ref{rem:S1acl} ; en
général, un pgcd constructif dans $k[a,b,X,Y,Z]$.

\section{Liste de formalisation}

\begin{enumerate}
\item Listes de Laurent sur un anneau, sens, produit de listes avec
  $\text{coef}(g\ell)=g\cdot\text{coef}(\ell)$ ; projection de degré ;
  lemmes~\ref{lem:regroupement} et~\ref{lem:denom}.
\item Modules à congruence : variantes de \texttt{FiniteCorrection},
  \texttt{ReesLifting}, \texttt{ACLZMT} ; \texttt{gen1} réduit aux
  $\beta_n$ ; \texttt{gen2} et \texttt{one-comb} modulo $N^0$, la
  restriction annulant $N^0$.
\item Le modèle et le lemme~\ref{lem:modele} ; la restriction.
\item Proposition~\ref{prop:fixe} sur les listes.
\item L'anneau $D=A[U,V]$ bigradué comme coefficients de \texttt{Laurent3}
  (déjà générique en $D$), $\phi$ et ses commutations ;
  proposition~\ref{prop:F1}, sous (E) en paramètre.
\item Proposition~\ref{prop:F2}, sous (S) et des générateurs de
  $\mathcal L$ en paramètres ; instance d'ACL ($\mathcal L=A^3$, (S0)
  trivial) avec (S1) par la remarque~\ref{rem:S1acl}.
\item (E) par bases de Gröbner (Coquand--Persson, théorème 11) et
  élimination ; (S1) en général par le pgcd ; noyaux de matrices.
\end{enumerate}
Les points 1 à 6 sont des récurrences sur des listes et des calculs
coefficient par coefficient du même type que ceux déjà vérifiés ; le
point 7 est le seul développement substantiel nouveau.

\end{document}
